%!TEX root = ../main.tex
% =============================================================================
%  六、问题 3：多阶段滚动调整模型      目标 3.5 页
%  评审在问：多阶段做对了吗？
%  本节含：① 后覆盖前的决策结构与非预期性的实现（式 (26) 场景树约束）
%          ② 分段线性凸偏差费用的精确线性化（凸性 → 无需 0-1 变量）
%          ③ 多阶段滚动模型与 MPC 实现（算法 1）
%          ④ 信息价值 VoI：直接回答"是否需要引入其他时刻预报"
%          ⑤ 结果与口径依赖
%  口径：C-09 在本节首次定义；表 tab:q3-epochs 给出四段槽位划分
% =============================================================================

\section{问题 3：多阶段滚动调整模型}\label{chap:q3}

\subsection{后覆盖前的决策结构}

问题 3 在问题 2 上增加两个要素：四个时刻 \(m\in\{0,6,12,18\}\) 均有未来 24~h 整点光伏预报，决策随之
\textbf{滚动修正}；并引入\textbf{两类不对称惩罚}（总费用含计划购电费、紧急购电费与调整相关费用）。
沿用 \chapref{chap:data} 的槽位读法：槽位 \(t=0,\dots,143\) 表示 \([10t,10(t+1))\) 分钟，
\(\Delta t=1/6\)~h。决策量为 \(0{:}00\) 制定的\textbf{计划购电量} \(E^{\text{plan}}_t\)（第一阶段承诺）、
滚动修正后的\textbf{调整购电量} \(E^{\text{adj}}_t\)、由实际负载与光伏决定的\textbf{紧急购电量} \(g_t\)，
以及充放电量 \(c_t,d_t\) 与时段末储电量 \(S_t\)。

四个预报时刻把一天切成四段，槽位 \(t\) 的最终执行量由\textbf{最后一次覆盖它的那次决策}给出
（表~\ref{tab:q3-epochs}）：\(0{:}00\) 用当次预报制定全天 144 槽计划，此后每个决策时刻只用\textbf{当次
预报}对\textbf{其后至 \(24{:}00\) 的剩余时段}重新决策，但\textbf{只执行到下一个决策时刻}。

\begin{table}[htbp]
  \centering
  \caption{四次预报的决策时刻、生效槽位与可用预报（时间范围：2025-02-01 至 12-31 的每一天；
            槽位单位：10 min 槽，共 144 槽；预报来源：附件 3）}
  \label{tab:q3-epochs}
  \footnotesize
  \setlength{\tabcolsep}{4pt}
  \begin{tabularx}{\linewidth}{C{1.6cm} L{3.0cm} C{2.5cm} C{2.4cm} Y}
    \toprule
    决策时刻 & 决策内容 & 生效槽位 & 生效时段 & 重解范围（只执行到下一时刻） \\
    \midrule
    \(0{:}00\) & 计划购电量 \(E^{\text{plan}}\)（144 槽） & \(0\)--\(35\) & \(0{:}00\)--\(6{:}00\) & \(0\)--\(143\) \\
    \(6{:}00\)  & 调整购电量 \(E^{\text{adj}}\) & \(36\)--\(71\) & \(6{:}00\)--\(12{:}00\) & \(36\)--\(143\) \\
    \(12{:}00\) & 调整购电量 \(E^{\text{adj}}\) & \(72\)--\(107\) & \(12{:}00\)--\(18{:}00\) & \(72\)--\(143\) \\
    \(18{:}00\) & 调整购电量 \(E^{\text{adj}}\) & \(108\)--\(143\) & \(18{:}00\)--\(24{:}00\) & \(108\)--\(143\) \\
    \bottomrule
  \end{tabularx}
\end{table}

\textbf{非预期性在实现层面自然满足。} 非预期性要求任一时刻的决策只依赖该时刻及之前已实现的信息；
本文实现有两个环节天然满足，无需显式施加式~\eqref{eq:q3-nonanticip}：其一，取预报值的函数只按
\textbf{提前期}索引——槽位 \(t\) 在时刻 \(m\) 取附件 3 中 \(m\) 时刻发布行的第
\(k=(h_t-m)\bmod24\) 列，故 \(m\) 时刻的子问题\textbf{在数据结构上取不到}其后发布的列；其二，
每个子问题只覆盖 \(t\ge s_m=6m\)，已执行槽位被冻结，决策不可能回头改写 \(t<s_m\) 的历史。
为可证伪地检验这一点，本文做\textbf{预报置换检验}：把某一次预报在 12 天之间随机置换后重解，
记录四个控制区内决策量的最大变化。结果是置换 \(0{:}00\) 预报会显著改变全天决策（说明检验有效力），
而\textbf{任何一次置换都不改变早于该预报发布时刻的槽位决策}（该区变化恰为 \(0.000\)~kWh）；
置换 \(18{:}00\) 预报时四区变化全为 \(0.000\)，这是数据事实而非检验失效——该次预报覆盖的槽
\(108\)--\(143\) 上几乎恒为 0（\(19{:}00\) 以后全年恒为 0），对晚间决策本无新信息
（逐区明细见 \appref{tab:na-app}{E 非预期性置换明细}）。

\subsection{偏差费用的凸性与精确线性化}

偏差费用 \(c^{\text{buy}}_t\) 关于 \(E^{\text{adj}}_t\) 是\textbf{凸的分段线性函数}：调增（\allowbreak
\(E^{\text{adj}}_t\ge E^{\text{plan}}_t\)）斜率为 \(1.5p_t\)，调减斜率为 \(0.5p_t\)，故

\begin{equation}
  c_t^{\text{buy}}
  =p_t\min\bigl(E^{\text{plan}}_t,E^{\text{adj}}_t\bigr)
  +0.5p_t\bigl(E^{\text{plan}}_t-E^{\text{adj}}_t\bigr)^{+}
  +1.5p_t\bigl(E^{\text{adj}}_t-E^{\text{plan}}_t\bigr)^{+}.
  \label{eq:q3-cost}
\end{equation}

凸性意味着\textbf{费用只惩罚偏离的方向}：调减按半价结算差额、调增按 1.5 倍结算。引入
\(u_t,v_t\ge0\) 表示超额量与违约量，

\begin{equation}
  E^{\text{adj}}_t-E^{\text{plan}}_t=u_t-v_t,\qquad
  c_t^{\text{buy}}=p_tE^{\text{plan}}_t+1.5p_tu_t-0.5p_tv_t ,
  \label{eq:q3-lin}
\end{equation}

后式由 \(p_t\min(E^{\text{plan}}_t,E^{\text{adj}}_t)=p_tE^{\text{plan}}_t-p_tv_t\) 得到。该线性化是
\textbf{精确的}，无需 0--1 变量：费用中与 \((u_t,v_t)\) 有关的部分为 \(1.5p_tu_t-0.5p_tv_t\)，
若 \(u_t,v_t\) 同时为正，两边同减 \(\min(u_t,v_t)\) 可使目标下降 \(p_t\min(u_t,v_t)>0\)，
故最优解必有 \(\min(u_t,v_t)=0\)。互斥性的关键是两条分段斜率跨越基准电价 \(1\times p_t\)
（\(0.5<1<1.5\)），使 \(u_t,v_t\) 的有效系数均为正，最优解只能落在折点上。
\textbf{适用边界}：若违约倍数取到 \(1\) 及以上（本题为 \(0.5\)），折点两侧斜率不再跨过基准电价，
互斥性不再自动成立，必须引入 0--1 变量。数值上，\(334\times144=48\,096\) 个“槽-日”中
\(\min(u_t,v_t)>10^{-6}\)~kWh 的违反数为 \textbf{0}。

\subsection{非预期性约束与多阶段滚动模型}\label{sec:q3-mpc}\label{sec:q3-roll}

以四个决策时刻为分支点构造\textbf{场景树}，记 \(\tau(t)\) 为不晚于槽位 \(t\) 的最后一个决策时刻。
\(0{:}00\) 的计划量对所有场景相同；槽位 \(t\) 的调整量只能依赖截至 \(\tau(t)\) 的信息，即

\begin{equation}
  E^{\text{adj}}_t(\omega)=E^{\text{adj}}_t(\omega'),
  \qquad
  \forall\,\omega,\omega'\in\Omega:\ \
  \omega|_{[0,\tau(t)]}=\omega'|_{[0,\tau(t)]},
  \label{eq:q3-nonanticip}
\end{equation}

其中 \(\omega|_{[0,\tau]}\) 表示场景 \(\omega\) 在时刻 \(\tau\) 之前的历史。式~\eqref{eq:q3-nonanticip}
展开成 LP，等价于按场景树中\textbf{深度为 \(\tau(t)\) 的节点}把场景分组、对同组场景施加变量相等的线性
约束。本文不显式展开场景树，而采用与它在确定性等价下一致的滚动实现（\S\ref{sec:q3-mpc}）：
变量规模与场景数解耦，非预期性由“数据可取范围 \(+\) 已执行槽位冻结”自动满足。

记 \(\Omega_m\) 为时刻 \(m\) 用于剩余时段的场景集（权重 \(w_\omega\)，按日总预报残差分 60 层、
层内取 medoid、权重取层样本占比，与 \ref{sec:q2-scen} 同一份存盘场景集），\(s_m=6m\)。综合
式~\eqref{eq:q3-cost}--\eqref{eq:q3-nonanticip}，问题 3 的多阶段模型为

\begin{subequations}\label{eq:q3}
\begin{align}
  \min\quad
  & \sum_{t=0}^{143}p_tE^{\text{plan}}_t
    +\sum_{\omega\in\Omega}w_\omega\sum_{t=0}^{143}
    \Bigl[1.5p_tu_t(\omega)-0.5p_tv_t(\omega)+5p_tg_t(\omega)\Bigr] \label{eq:q3-obj}\\
  \text{s.t.}\quad
  & E^{\text{adj}}_t(\omega)-E^{\text{plan}}_t=u_t(\omega)-v_t(\omega),\quad u_t(\omega),v_t(\omega)\ge0, && \forall t,\omega \label{eq:q3-diff}\\
  & E^{\text{plan}}_t-c_t+d_t\ \ge\ \bigl(L_t-\widehat{PV}^{(m)}_t\bigr)\Delta t, && \forall t \label{eq:q3-planfeas}\\
  & g_t(\omega)\ \ge\ \bigl(L_{\omega,t}-PV_{\omega,t}\bigr)\Delta t+c_t-d_t-E^{\text{adj}}_t(\omega),\quad g_t(\omega)\ge0, && \forall t,\omega \label{eq:q3-gap}\\
  & S_t=S_{t-1}+\eta_cc_t-d_t/\eta_d,\quad \underline S\le S_t\le\overline S,\quad 0\le c_t,d_t\le\overline P\Delta t,\quad S_0=S^{init},\ S_{143}\ge S_0, && \forall t \label{eq:q3-soc}\\
  & E^{\text{plan}}_t\ \text{与}\ \omega\ \text{无关};\qquad E^{\text{adj}}_t\ \text{满足式}~\eqref{eq:q3-nonanticip}. && \label{eq:q3-na-ref}
\end{align}
\end{subequations}

式~\eqref{eq:q3-gap} 的缺口以\textbf{调整后}的 \(E^{\text{adj}}_t(\omega)\) 为基准（假设 A7）；
终端条件 \(S_{143}\ge S_0\) 防止滚动优化在窗口末端放空储能换取短期收益（假设 A6）。参数取
\(\eta_c=\eta_d=0.9\)、\(\overline P\Delta t=833.3333\)~kWh、\([\underline S,\overline S]=[1200,10\,800]\)~kWh、
\(S^{init}=6000\)~kWh。

\textbf{MPC 实现。} 式~\eqref{eq:q3} 的场景树 LP 按决策时刻分解为四个子问题：在时刻 \(m\) 以当前
储电量 \(S_{s_m-1}\) 为初值、当日初值 \(S_0\) 为终端下界，对 \(t\ge s_m\) 求解

\begin{equation}
  \min\ \sum_{t=s_m}^{143}
  \Bigl[(1-\delta_m)p_tE_t+\delta_m\bigl(1.5p_tu_t-0.5p_tv_t\bigr)\Bigr]
  +\sum_{\omega\in\Omega_m}w_\omega\sum_{t=s_m}^{143}5p_tg_{t,\omega},
  \qquad
  \delta_m=\begin{cases}0,&m=0\\ 1,&m>0\end{cases}
  \label{eq:q3-mpc}
\end{equation}

约束为式~\eqref{eq:q3-diff}--\eqref{eq:q3-soc} 在剩余槽位上的限制；\(m>0\) 时计划项已是常数，
故从目标中略去。式~\eqref{eq:q3-mpc} 把“重解范围”与“执行范围”分离：求解后\textbf{只把}
\([s_m,s_{m+1})\) 上的决策写入实际轨迹，其余解丢弃，待下一个预报时刻用更新的预报重优化——充放电量
因此随信息滚动修正，而非冻结在 \(0{:}00\)。全部物理约束由结果文件独立复算：供给缺口残差
\(\le3.5\times10^{-13}\)~kWh、储电量递推残差 \(\le1.1\times10^{-11}\)~kWh、储电量与功率越界数均为 0；
每个结果日重解 4 次子问题（共 1336 次，另加 1 月预热期），每次均为 LP 并由 HiGHS 精确求解。
完整流程见算法~\ref{alg:rolling}。

\begin{algorithm}[htbp]
  \caption{微网日前计划——多阶段滚动调整的一体化求解流程}
  \label{alg:rolling}
  \small
  \begin{algorithmic}[1]
    \Require 电价 \(p_t\)、负载 \(L_t\)、实际光伏 \(P^{act}_t\)、各发布时刻预报 \(P^{fe}_t\)
             (\(e\in\{0,6,12,18\}\))、储能参数 \(\eta_c,\eta_d,\overline P,\underline S,\overline S\)、
             初值 \(S_0=6000\)~kWh
    \Ensure 计划购电量 \(E^{\text{plan}}\)、调整购电量 \(E^{\text{adj}}\)、实际充放电量 \(c,d\)、
            紧急购电量 \(g\)、储电量轨迹 \(S\)
    \For{每一天 \(d\)}
      \State \textbf{计划层}：以 \(0{:}00\) 预报为光伏输入求解式~\eqref{eq:q2-obj}，得
             \(E^{\text{plan}},c,d,S\)（问题 2/4-2 到此为止）
      \For{每个调整时刻 \(e\in\{6,12,18\}\)}
        \State \textbf{调整层}：以该时刻预报重解\textbf{剩余时段} \([6e/6,\,144)\) 的 LP，
               目标含偏差惩罚与 \(5p_t\) 紧急购电，约束加\textbf{非预期性}
               （槽位 \(<6e/6\) 的决策冻结）
        \State 后覆盖前：该次解覆盖其后所有槽位的执行值
      \EndFor
      \State \textbf{结算}：按附件 2 实际负载与光伏计算缺口
             \(g_t=\bigl[(L_t-P^{act}_t)\Delta t+c_t-d_t-E^{\text{adj}}_t\bigr]^{+}\)，
             紧急购电费 \(5\sum_tp_tg_t\)
      \State \textbf{一致性校验}：核验电量平衡、储电量递推、区间与功率界、
             充放互斥与式~\eqref{eq:q1-energy} 的能量恒等式
      \State \textbf{递推}：\(S_0^{(d+1)}\leftarrow S_{144}^{(d)}\)
    \EndFor
    \State \Return 全部策略与费用构成
  \end{algorithmic}
\end{algorithm}

\subsection{信息价值：是否需要引入其他时刻的预报}\label{sec:voi}

题面直接追问该问题，故须给出\textbf{量化答案}与\textbf{成本侧}。设计\textbf{预报时刻数对照实验}：
\(n\in\{1,2,4\}\) 三组使用\textbf{同一场景集、同一实际轨迹回放、同一费用口径}（\calref{C-09}），
只有“计划时能用到几次预报”不同——\(n=1\) 只用 \(0{:}00\) 预报（即问题 2 的结构）、\(n=2\) 用
\(0{:}00\) 与 \(12{:}00\)、\(n=4\) 为题面给定的四次预报。附件 3 只提供这四个发布时刻，故
\(n=8,24\) 在本数据下\textbf{并不存在}：任何加密实验都须自行合成预报，其结论取决于合成规则而非
数据。本文不做这种合成，改用\textbf{完全信息上界}：\(0{:}00\) 计划仍按 \(0{:}00\) 预报制定，
但从 \(6{:}00\) 起的调整假定已知当日实际净需求，其费用\textbf{不高于}任何有限预报时刻数配置，
是“还能再省多少”的\textbf{与数据无关的上界}。记

\begin{equation}
  \mathrm{VoI}(n)=\frac{C(1)-C(n)}{C(1)}\times100\%,
  \label{eq:q3-voi}
\end{equation}

其中 \(C(n)\) 为可用 \(n\) 个预报时刻时、同一场景集同一实际轨迹下的总费用。结果见表~\ref{tab:voi}

\begin{table}[htbp]
  \centering
  \caption{预报时刻数对照实验：总费用、边际节省与信息价值
           （同一场景集、同一实际轨迹、同一偏差量定价\calref{C-09}；
            时间范围：2025-02-01 至 12-31 共 334 天；样本：60 个分层 medoid 代表场景 \ensuremath{\times} 4 组配置；
            单位：费用 元，电量 kWh）}
  \label{tab:voi}
  \footnotesize
  \setlength{\tabcolsep}{4pt}
  \begin{tabularx}{\linewidth}{L{3.0cm} C{2.1cm} C{1.9cm} C{1.6cm} C{1.8cm} Y}
    \toprule
    可用预报时刻 & 总费用/\unit{元} & 边际节省/\unit{元} & 累计 \(\mathrm{VoI}\) & 已实现占比 & 紧急购电量/\unit{kWh} \\
    \midrule
    仅 \(0{:}00\)（\(n=1\)） & 14\,264\,792.13 & --- & 0 & 0 & 172\,015.82 \\
    \(0{:}00\)、\(12{:}00\)（\(n=2\)） & 14\,131\,564.16 & 133\,227.98 & 0.93\% & 34.1\% & 122\,524.83 \\
    四时刻（\(n=4\)，题面给定） & 13\,932\,402.95 & 199\,161.20 & 2.33\% & 85.1\% & 64\,522.09 \\
    完全信息（可获信息上界） & 13\,874\,340.39 & 58\,062.56 & 2.74\% & 100\% & 6\,134.76 \\
    \bottomrule
  \end{tabularx}
  \\[1ex]
  {\footnotesize 表注：“已实现占比”的分母 \(C(1)-C_{\mathrm{PI}}=390\,451.74\) 元是完全信息下最多能省的
  金额。}
\end{table}

结论有三条。① \textbf{滚动调整有效}：总费用降 \(2.33\%\)、紧急购电量降 \(62.49\%\)。
② \textbf{收益边际递减}：\(1\to2\) 次加密省 \(133\,227.98\) 元、\(2\to4\) 次省 \(199\,161.20\) 元，
\(n=4\) 已实现可获得信息价值的 \(\mathbf{85.1\%}\)；\textbf{再增加预报最多只能再省}
\(58\,062.56\) 元（占 \(n=4\) 总费用的 \(0.42\%\)）。③ \textbf{边际增益主要来自 \(6{:}00\)
对上午的更新}：\(199\,161.20\) 元不是三次预报的平均贡献，由置换检验，\(18{:}00\) 预报在其控制的槽上
几乎恒为 0、置换它不改变任何决策，而上午正是 \(0{:}00\) 预报系统性偏离最大（\chapref{chap:data}）
与紧急购电最集中的时段。

\textbf{成本侧与结论。} 把收益上界 \(58\,062.56\) 元摊到 334 天，等价于每天新增预报投入不超过
\(173.84\) 元。故本文的回答是：\textbf{四次预报已经足够，在预报获取成本可以忽略之前不建议继续加密
预报时刻}；若要追加投入，应优先投向\textbf{提高现有四次预报的精度}——相同时刻数下剩余误差仍对应
\(58\,062.56\) 元（\(0.42\%\)），只能靠精度改进消除 \evid{C-Q3-01}。适用边界：① 只在本题惩罚
倍数（违约 \(0.5p_t\)、超额 \(1.5p_t\)、紧急 \(5p_t\)）下成立；② 只在附件 3 的预报质量下成立，
预报误差成倍减小会使边际收益更早饱和；③ \(n=8,24\) 为数据中不存在的配置，以其上界代替，
故“不建议加密”是对\textbf{收益侧上界}的判断。实现正确性自检：四组费用满足不等式链

\begin{equation}
  C_{\mathrm{PI}}\le C_{n=4}\le C_{n=1}:\qquad
  13\,874\,340.39\le13\,932\,402.95\le14\,264\,792.13\quad(\text{元});
  \label{eq:q3-chain}
\end{equation}

若数值违反该链，说明非预期性、信息集或结算口径实现有误。

\subsection{求解结果与口径依赖}

在同一场景集上求解式~\eqref{eq:q3}，四项费用为：计划相关 \(13\,216\,203.48\) 元、违约
\(210\,674.71\) 元、超额 \(262\,628.79\) 元、紧急购电 \(242\,895.98\) 元，合计
\(13\,932\,402.95\) 元（1\,393.24 万元），相对无储能基线 \(16\,407\,319.63\) 元节省 \(15.08\%\)；
紧急购电量 \(64\,522.09\)~kWh，出现在 302/334 天（\calref{C-09}）。

\textbf{调整机制省在哪里。} 与“只按 \(0{:}00\) 计划结算”的问题 2 结构（计划购电费
\(13\,637\,552.90\) 元 \(+\) 紧急购电费 \(627\,239.24\) 元）相比，调整机制把\textbf{紧急购电费压低
\(384\,343.26\) 元}，代价是计划相关费用净增 \(51\,954.08\) 元（计划项减少 \(421\,349.41\) 元、
违约与超额合计增加 \(473\,303.50\) 元），相抵净省 \(332\,389.18\) 元。\(0{:}00\) 计划购电量合计
\(22\,220\,333.20\)~kWh，调整后 \(21\,843\,259.36\)~kWh（净减 \(377\,073.84\)~kWh），其中调减
\(645\,614.42\)~kWh、调增 \(268\,540.58\)~kWh，二者之差恰为净减少量（口径自检）。\(0{:}00\) 计划量
在四种信息结构下\textbf{完全相同}，信息价值的全部来源都在\textbf{调整}环节。

\textbf{与问题 2 的对照与口径依赖。} 两者同为附件 1 固定电价、同一场景集、同一实际轨迹与同一费用
口径，唯一差别是三次滚动预报是否可用；费用差额即信息价值 \(332\,389.18\) 元（\(2.33\%\)），相对
无储能基线的节省率由 \(13.06\%\) 提高到 \(15.08\%\)。该降幅\textbf{依赖于 \calref{C-09} 的偏差量定价
口径}：在“计划购电量全额结算 \(+\) 差额惩罚”的读法下，计划项不再随调整量下降，净收益会减少
\(\sum_tp_tv_t=421\,349.41\) 元并转为负值（\chapref{chap:verify} 给出对照），故该答案与市场规则
（惩罚倍数与结算方式）绑定，不能脱离口径单独引用。按题面表 1／表 2／表 3 格式给出的指定日期结果见
\chapref{chap:results}表~\ref{tab:key4}，其中“调整购电量”一列填的是\textbf{最后一次覆盖该槽位的决策}
给出的购电量（表~\ref{tab:q3-epochs}），与 \(0{:}00\) 计划量之差即该槽的偏差量。
